\documentclass[11pt]{article}
\usepackage{fontspec}
\usepackage[french]{babel}
\usepackage{amsmath}
\usepackage{amssymb}
\usepackage{amsthm}
\usepackage{mathrsfs}
\usepackage{tikz-cd}
\usepackage{enumitem}
\usepackage{lmodern}

\newcommand{\printedpage}[1]{\marginpar{\footnotesize p.~#1}}
\newcommand{\msp}[1]{\marginpar{\footnotesize [#1]}}
\newcommand{\Ob}{\operatorname{Ob}}
\newcommand{\Filt}{\mathrm{Filt}}
\newcommand{\Ind}{\operatorname{Ind}}
\newcommand{\card}{\operatorname{card}}
\newcommand{\cardred}{\operatorname{card\,r\acute{e}d}}

\begin{document}

\noindent\printedpage{59}% head: A. Grothendieck : Dérivateurs, Ch. XVIII, Édition des Universités de Montpellier II et Paris VII
% foot: Version 17 décembre 2025
sauf qu'il faut faire attention à la majoration de cardinaux, $\card I \leqslant \pi'^{\pi}$. Mais si $F \in \Ob\mathcal{F}^{\pi'}$, comme par (4.7.6.1)
\[
\mathcal{F}^{\pi'} \;\subset\; \underbrace{\mathcal{M}_{\pi'}}_{\subset\mathcal{M}_{\pi'^{\pi}}\subset\Filt'^{\pi'^{\pi}}} \;\subset\; \Filt'^{\pi'^{\pi_0}} ,
\]
donc a fortiori
\[
\mathcal{F}^{\pi'} \;\subset\; \Filt^{\pi'^{\pi}} \underbrace{\;=\;}_{\text{si } \pi'^{\pi}\geqslant\pi_1} \Filt^{\pi'^{\pi}} , % [sic] first \Filt^{\pi'^{\pi}} -> \Filt'^{\pi'^{\pi}}, otherwise the equality under the brace would be trivial (cf. the preceding display)
\]
d'où, par 4.7.3, avec $\pi'$ remplacé par $\pi'^{\pi}=\pi''$, la possibilité de prendre $I$ avec
\[
\card I \;\leqslant\; \pi''^{\pi} \;=\; (\pi'^{\pi})^{\pi} \;=\; \pi'^{\pi\cdot\pi} \;=\; \pi'^{\pi} ,
\]
ce qui achève la démonstration.

\noindent\msp{page 175}\textbf{Remarque 4.7.7.2.} J'ai été obligé, dans l'énoncé de la condition c), de faire intervenir un deuxième cardinal $\pi_1$, et ceci faut de savoir si, pour $\pi''$ tel que $\pi''=\pi''^{\pi_0}$, $\Filt^{\pi''}$ est stable par facteurs directs. (Il faudrait, pour bien faire, soit prouver qu'il est stable, donc qu'on a dans 4.7.1 c $\Longrightarrow$ b, i.e. $\mathcal{M}_\pi\subset\Filt^{\pi}$ (et non seulement c $\Longrightarrow$ b$'$, i.e. $\mathcal{M}_\pi\subset\Filt'^{\pi}$) si $\pi=\pi^{\pi_0}$, ou bien faire un contre-exemple!) Cet ingrédient un peu étrange, on sera hélas obligé de le traîner dans la définition des filtrations cardinales généralisées, si on ne veut s'interdire de regarder dans le cas présent $(\mathcal{F}^{\pi})_{\pi\geqslant\pi_0}$ comme une telle filtration, et ne s'autoriser que de regarder $(\mathcal{F}^{\pi})_{\pi\geqslant\pi_1}$ comme telle, ce qui me paraîtrait ridicule. Pratiquement (même si théoriquement l'introduction de $\pi_1$ dans la description d'une filtration cardinale généralisée était bel et bien nécessaire) on n'aura guère des cas où il faille faire mention de $\pi_1$ dans la condition c. En effet, $\pi\geqslant\pi_0$ et $\pi^{\pi_0}=\pi$ impliquent $\pi\geqslant 2^{\pi_0}$, et dès lors $\pi'^{\pi}\geqslant 2^{(2^{\pi_0})}$. Si donc $\pi_1\leqslant 2^{(2^{\pi_0})}$, i.e. $\card C_0\leqslant 2^{(2^{\pi_0})}$, alors la condition $\pi^{\pi'}\geqslant\pi_1$ dans c) est automatiquement satisfaite, et il n'y a donc pas à la mentionner. Or il est bien rare qu'on ait à faire à des petites catégories $C_0$, dont le cardinal soit $>$ puissance % [sic] faut -> faute, "faute de savoir" intended; [sic] \pi^{\pi'}\geqslant\pi_1 -> \pi'^{\pi}\geqslant\pi_1, the condition is written \pi'^{\pi}\geqslant\pi_1 elsewhere (cf. 4.7.8)
\msp{page 176}
du continu $c$, et a fortiori, telle que $\card C_0 > 2^{c}$. Si par contre $\card C_0\leqslant 2^{c}$, comme $2^{c}\leqslant 2^{(2^{\pi_0})}$ (puisque $\pi_0\geqslant\aleph_0$), on aura bien $\pi_1\leqslant 2^{(2^{\pi_0})}$, et on est tranquille.

\textbf{Définition 4.7.8.} Soit $\mathcal{M}$ une $\mathfrak{U}$-catégorie. On appelle $\mathfrak{U}$-\textit{filtration cardinale généralisée} (ou simplement filtration cardinale généralisée, quand $\mathfrak{U}$ est sous-entendu) la donnée d'un cardinal infini $\pi_0$, et d'une filtration $(\mathcal{F}^{\pi})_{\pi\geqslant\pi_0}$ de $\mathcal{M}$ par des sous-catégories strictement pleines $\mathcal{F}^{\pi}$, indexées par les $\mathfrak{U}$-cardinaux $\geqslant\pi_0$, et satisfaisant les conditions a) b) c) énoncées dans 4.7.7, où dans c) on pose
\[
\pi_1 \;=\; \sup(\pi_0,\cardred\mathcal{F}^{\pi_0}) \, . \tag*{(4.7.8.1)}
\]
(Donc si $\cardred\mathcal{F}^{\pi_0}\leqslant 2^{(2^{\pi_0})}$, il n'y a pas à tenir compte de la condition $\pi'^{\pi}\geqslant\pi_1$, cf. 4.7.7.1.) On dit que la filtration cardinale généralisée (fcg) est \textit{accessible}, si elle satisfait de plus la condition d). On dit qu'elle est \textit{canonique} si on a même
\[
\mathcal{F}^{\pi} \;=\; \mathcal{M}_{\pi} \tag*{(4.7.8.2)}
\]
pour tout $\pi\geqslant\pi_0$.

\textbf{Proposition 4.7.9.} \textit{Soient $(\mathcal{F}^{\pi})_{\pi\geqslant\pi_0}$ une fcg sur la $\mathfrak{U}$-catégorie $\mathcal{M}$, et $\pi>\pi_0$ un cardinal tel que $\pi^{\pi_0}=\pi$. Alors on a}
\printedpage{60}% head: A. Grothendieck : Dérivateurs, Ch. XVIII, Édition des Universités de Montpellier II et Paris VII
% foot: Version 17 décembre 2025
\msp{page 177}
\[
\left\{
\begin{aligned}
&\mathcal{M}_\pi\subset\mathcal{F}^{\pi}\quad\text{\textit{(si $\pi=\pi^{\pi_0}$),}}\\
&\text{\textit{donc} }\mathcal{M}_\pi=\mathcal{F}^{\pi}\text{ \textit{si de plus la fcg est} accessible.}
\end{aligned}
\right.
\tag*{(4.7.9.1)}
\]

\textsc{Démonstration} (\footnote[75]{\texttt{[Dans ce qui suit, $\mathcal{F}^{\pi}$ a parfois été écrit $\mathcal{F}_{\pi}$, souvent corrigé en $\mathcal{F}^{\pi}$ par l'auteur. Nous écrivons $\mathcal{F}^{\pi}$ partout.]}}). Appliquant la condition c) des fcg à un $F\in\Ob\mathcal{M}_\pi$, on trouve que $F$ est facteur direct d'un $x_i\in\Ob\mathcal{F}^{\pi}$. Comme $\mathcal{F}^{\pi}$ est stable par $\varinjlim$ grandes devant $\pi_0$ et de cardinal $\leqslant\pi$, et que \textit{de plus} $\pi>\pi_0$, on conclut que :

\textbf{Lemme 4.7.9.1} \texttt{[deux fois 4.7.9.1]}. \textit{Si $\pi>\pi_0$ (\footnote[76]{Mais $\mathcal{F}^{\pi_0}$ n'est pas forcément stable par facteurs directs, comme on voit en prenant $\mathcal{M}=\Ind_{\pi_0}(C)$, $C$ une petite catégorie non karoubienne, et la filtration cardinale pour les $\Filt'^{\pi}$ -- on note que $\Filt'^{\pi_0}=\Filt^{\pi_0}\simeq e$, vu qu'une catégorie grande devant $\pi_0$ et de cardinal $\leqslant\pi_0$ est cofinale à $e$.}), $\mathcal{F}^{\pi}$ est stable par facteurs directs.}

D'où résulte que $F\in\mathcal{F}^{\pi}$, d'où (4.7.9.1).

\textbf{Corollaire 4.7.10.} \textit{Supposons la fcg accessible. Pour tout cardinal $\pi\geqslant\pi_0$, on a}
\[
\mathcal{F}^{\pi} \;\subset\; \mathcal{M}_\pi \;\subset\; \mathcal{F}^{\pi^{\pi_0}} \; ( \;=\; \mathcal{M}_{\pi^{\pi_0}}) \, . \tag*{(4.7.10.1)}
\]

En effet, $\mathcal{F}^{\pi}\subset\mathcal{M}_\pi$ par définition de fcg accessibles, et $\mathcal{M}_\pi=\mathcal{M}_{\pi^{\pi_0}}$ (puisque $\pi\leqslant\pi^{\pi_0}$), or $\mathcal{M}_{\pi^{\pi_0}}=\mathcal{F}^{(\pi^{\pi_0})}$ en vertu de (4.7.9.1) appliqué à $\pi'=\pi^{\pi_0}$ au lieu de $\pi$ (vu que $\pi'^{\pi_0}=\pi'$), d'où (4.7.10.1). On notera que cette inclusion généralise (4.7.9.1), elle donne $\mathcal{F}^{\pi}=\mathcal{M}_\pi$ si $\pi=\pi^{\pi_0}$.

\textbf{4.7.11.} Il me semble que les seules filtrations cardinales généralisées utiles sont celles qui sont \textit{accessibles.} Par la suite, je vais me borner à ce cas. Comme alors $\mathcal{F}^{\pi}\subset\mathcal{M}_\pi$, on en conclut
\msp{page 178}
que le foncteur canonique
\[
\Ind_\pi(\mathcal{F}^{\pi}) \;\hookrightarrow\; \mathcal{M} \tag*{(4.7.11.1)}
\]
prolongeant l'inclusion $\mathcal{F}^{\pi}\subset\mathcal{M}$, de façon à commuter aux $\varinjlim$ grandes devant $\pi$, est \textit{pleinement fidèle.} De plus, lorsque $\pi^{\pi_0}=\pi$, la condition c) de fcg nous dit entre autres que ce foncteur est essentiellement surjectif, donc que c'est une \textit{équivalence de catégories} :
\[
\Ind_\pi(\mathcal{F}^{\pi}) \xrightarrow{\ \sim\ } \mathcal{M} \qquad (\text{si } \pi=\pi^{\pi_0}). \tag*{(4.7.11.2)}
\]

\textbf{Proposition 4.7.12.} \textit{Soit $(\mathcal{F}^{\pi})_{\pi\geqslant\pi_0}$ une fcg accessible (fcga) sur la $\mathfrak{U}$-catégorie $\mathcal{M}$, et soit $\pi$ un cardinal $\geqslant\pi_0$. Les conditions suivantes sont équivalentes.}
\begin{enumerate}[label=\textit{\alph*)}]
\item \textit{Le foncteur pleinement fidèle (4.7.11.1) est essentiellement surjectif, i.e. est une équivalence de catégories.}
\item \textit{Pour tout $F\in\Ob\mathcal{M}$, $\mathcal{F}^{\pi}/F$ est une catégorie filtrante et grande devant $\pi$.}
\end{enumerate}

\noindent\printedpage{61}% head: A. Grothendieck : Dérivateurs, Ch. XVIII, Édition des Universités de Montpellier II et Paris VII
% foot: Version 17 décembre 2025
\textit{Quand ces conditions sont satisfaites, et si $\pi>\pi_0$, on a}
\[
\mathcal{F}^{\pi} \;=\; \mathcal{M}_\pi \, . \tag*{(4.7.12.1)}
\]

\textsc{Démonstration}. L'équivalence de a) et b) résulte aussitôt de 4.4.2 et de la définition des catégories grandes devant $\pi$ (p.~81). Que ces
\msp{page 179}
conditions impliquent $\mathcal{F}^{\pi}=\mathcal{M}_\pi$, i.e. $\mathcal{M}_\pi\subset\mathcal{F}^{\pi}$ dans le cas où $\pi>\pi_0$, se voit comme (4.7.9.1) en utilisant le lemme 4.7.9.1.

\textbf{Définition 4.7.13.}
\begin{enumerate}[label={$\arabic*^{\circ}$)}]
\item Un cardinal $\pi\geqslant\pi_0$ est dit \textit{cardinal utile} pour la fcga $(\mathcal{F}^{\pi})_{\pi\geqslant\pi_0}$, s'il satisfait les conditions équivalentes a) et b) de 4.7.12, et si de plus dans le cas où $\pi=\pi_0$, on a $\mathcal{F}^{\pi}=\mathcal{M}_\pi$. (Donc si $\pi$ est un cardinal utile, on a toujours $\mathcal{F}^{\pi}=\mathcal{M}_\pi$.) Dans le cas contraire, on dit que $\pi$ est un \textit{cardinal superflu} (pour le fcga donnée). % [sic] le fcga donnée -> la fcga donnée, gender agreement (feminine)
\item Un cardinal $\pi$ est dit \textit{régulier} pour la fcg donnée, s'il est utile, et si de plus la majoration de cardinaux énoncée dans (4.7.7.1) est valable.
\end{enumerate}

\textbf{4.7.14} Ainsi, tous les cardinaux $\pi\geqslant\pi_0$ tels que $\pi^{\pi_0}=\pi$ sont utiles, et même réguliers pour la filtration cardinale \texttt{[généralisée]} donnée. Notons que si tous les cardinaux $\pi\geqslant\pi_0$ sont utiles, alors on a $\mathcal{F}^{\pi}=\mathcal{M}_\pi$ pour tout $\pi\geqslant\pi_0$, donc la fcga est canonique. Notons aussi que dans tous les cas, si $(\mathcal{F}^{\pi})_{\pi\geqslant\pi_0}$ est une fcga, alors $(\mathcal{M}_\pi)_{\pi\geqslant\pi_0}$ en est une aussi.

\noindent\msp{page 180}En effet, la condition a) des fcg (cf. 4.7.7) résulte de la même condition pour $(\mathcal{F}^{\pi})$, et de (4.7.10.1), $\mathcal{M}_\pi\subset\mathcal{F}^{\pi^{\pi_0}}$. La condition b) est une propriété générale de la $\pi$-accessibilité, compte tenu que $(\mathcal{F}^{\pi_0})_0$ est une petite sous-catégorie génératrice formée d'objets $\pi_0$-accessibles de $\mathcal{M}$. La condition d) est satisfaite tautologiquement, et il reste à vérifier la condition c). Mais on a vu que pour $\pi=\pi^{\pi_0}$, on a $\mathcal{M}_\pi=\mathcal{F}^{\pi}$ (4.7.9.1), donc la condition c) pour $(\mathcal{M}_\pi)$ résulte de la même condition pour $(\mathcal{F}^{\pi})$. Ainsi, on trouve :

\textbf{Proposition 4.7.15.} \textit{Si $(\mathcal{F}^{\pi})_{\pi\geqslant\pi_0}$ est une fcga sur la $\mathfrak{U}$-catégorie $\mathcal{M}$, il en est de même de $(\mathcal{M}_\pi)_{\pi\geqslant\pi_0}$. De plus, tout cardinal utile (resp. régulier) pour la première filtration, l'est pour la seconde (appelée la fcg} canonique associée \textit{à la filtration cardinale} \texttt{[généralisée]} \textit{donnée).}

La dernière assertion résulte du fait que si $\pi$ est utile pour $(\mathcal{F}^{\pi})$, alors on a $\mathcal{F}^{\pi}=\mathcal{M}_\pi$.

\textbf{4.7.16.} Il me semble qu'on peut dire que la fcg canonique associée à une fcga donnée
\msp{page 181}
est toujours ``meilleure'' que celle-ci (si celle-ci n'est canonique, i.e. si les deux ne sont égales), à cause des deux faits suivants :
\begin{enumerate}[label=\alph*)]
\item justement parce qu'elle est \textit{canonique}, i.e. ne dépend d'aucun autre choix que celui du seul cardinal $\pi_0$, et
\item parce qu'il y a a priori plus des cardinaux utiles et de cardinaux réguliers pour la fcg canonique.
\end{enumerate}
Donc au total, il me semble que les filtrations cardinales \texttt{[généralisées]} les seules utiles sont celles qui sont canoniques.

\noindent\printedpage{62}% head: A. Grothendieck : Dérivateurs, Ch. XVIII, Édition des Universités de Montpellier II et Paris VII
% foot: Version 17 décembre 2025
Notons qu'une filtration cardinale (au sens de SGA 4 I 9, que je conserve ici) qui est de plus \textit{accessible} (i.e. $\mathcal{F}^{\pi}\subset\mathcal{M}_\pi$ pour tout $\pi\geqslant\pi_0$), n'est autre qu'une fcga canonique, telle que \textit{tous} les cardinaux sont réguliers (a fortiori utiles) pour ladite filtration. Je présume (sans avoir cependant construit d'exemple) que les filtrations cardinales \texttt{[généralisées]} canoniques des catégories de la forme $\mathcal{M}=\Ind_{\pi_0}(C)$, $C$ une petite catégorie \textit{karoubienne} (pour que $\pi_0$ soit un cardinal utile) ne sont pas forcément des fc (non généralisées, i.e. au sens SGA), i.e. je présume qu'il arrive qu'il y ait des cardinaux $\pi\geqslant\pi_0$ non réguliers, et même superflus.

\noindent\msp{page 182}Il me faut enfin étudier ce qui se passe quand on change le cardinal initial $\pi_0$ d'une fcga.

\textbf{Proposition 4.7.17.} \textit{Soit $(\mathcal{F}^{\pi})_{\pi\geqslant\pi_0}$ une fcga sur la $\mathfrak{U}$-catégorie $\mathcal{M}$, et soit $\pi'_0$ un cardinal $\geqslant\pi_0$. Alors $(\mathcal{F}^{\pi})_{\pi\geqslant\pi'_0}$ est une fcga sur $\mathcal{M}$, qui est canonique si celle de départ est canonique. Pour un cardinal $\pi\geqslant\pi'_0$, celui-ci est utile pour $(\mathcal{F}^{\pi})_{\pi\geqslant\pi'_0}$ si et seulement s'il l'est pour $(\mathcal{F}^{\pi})_{\pi\geqslant\pi_0}$. Pour qu'il soit régulier pour $(\mathcal{F}^{\pi})_{\pi\geqslant\pi'_0}$, il} suffit \textit{qu'il le soit pour $(\mathcal{F}^{\pi})_{\pi\geqslant\pi_0}$, et c'est aussi nécessaire quand on suppose de plus que $2^{\pi}\geqslant\pi'_1\overset{\text{déf}}{=}\sup(\pi'_0,\cardred\mathcal{F}^{\pi'_0})$, a fortiori quand $\pi\geqslant\pi'_1$.}

\textsc{Démonstration}. Les conditions a) b) sont trivialement satisfaites par la deuxième filtration, étant satisfaites par la première, ainsi que d) si elle l'est par la première. Il en est de même de la condition c) si $\pi=\pi^{\pi'_0}$, qui implique en effet $\pi^{\pi_0}=\pi$, puisqu'en général
\[
\pi \;\leqslant\; \pi^{\pi_0} \;\leqslant\; \pi^{\pi'_0} \, .
\]
Reste à examiner la majoration de cardinaux
\[
\card I \;\longrightarrow\; \pi'^{\pi} \qquad \text{si } F\in\Ob\mathcal{F}^{\pi'} \ \underline{\text{et}}\ \pi'^{\pi}\geqslant\pi'_1 \, , % [sic] \longrightarrow -> \leqslant, a majoration (inequality) of cardinals is meant
\]
où cette fois $\pi'_1=\sup(\pi'_0,\cardred\mathcal{F}^{\pi'_0})$. En tout état de cause, on a $\pi'_1\geqslant\pi_1$, donc
\msp{page 183}
$\pi'^{\pi}\geqslant\pi'_1$ implique $\pi'^{\pi}\geqslant\pi_1$, donc on peut appliquer la condition c) des fcg pour $(\mathcal{F}^{\pi})_{\pi\geqslant\pi_0}$, pour conclure qu'elle reste encore vérifiée. Donc en somme, \textit{toutes} les conditions des fcg (resp. fcga) sont trivialement vérifiées pour la fcg \textit{restreinte} de la première. De même, le fait qu'un cardinal $\pi'\geqslant\pi'_0$ est utile est trivialement indépendant de quelle des deux filtrations on considère, elle ne dépend que de la donnée du cardinal $\pi'$ et de la sous-catégorie $\mathcal{F}^{\pi'}$ de $\mathcal{M}$. Reste la majoration de cardinaux, dans le cas où on s'intéresse à la \textit{régularité} de $\pi'$ pour les fcga envisagées. La majoration exigée est la même, mais la condition préalable de validité de cette majoration, savoir
\[
\pi'^{\pi} \;\geqslant\; \pi'_1 \qquad \text{resp.} \qquad \pi'^{\pi} \;\geqslant\; \pi_1 \, ,
\]
n'est pas la même. Comme $\pi'_1\geqslant\pi_1$, on voit que la condition préalable est moins restrictive dans le cas de la première filtration, que de la seconde, d'où le ``il suffit'' dans la dernière assertion. Mais lorsque $2^{\pi}\geqslant\pi'_1$, même la plus restrictive est automatique\msp{page 184}ment vérifiée, puisque $\pi'^{\pi}\geqslant 2^{\pi}\geqslant\pi'_1$, et alors la régularité de $\pi$ pour la deuxième filtration implique la régularité pour le première, q.e.d. % [sic] le première -> la première, gender agreement (feminine)

\noindent\printedpage{63}% head: A. Grothendieck : Dérivateurs, Ch. XVIII, Édition des Universités de Montpellier II et Paris VII
% foot: Version 17 décembre 2025
\textbf{Corollaire 4.7.18.} \textit{Supposons $\pi_0$ régulier pour la fcga (\footnote[77]{Il suffit pour cela que $\pi_0$ soit utile, i.e. $\mathcal{M}\xleftarrow{\ \sim\ }\Ind_{\pi_0}(\mathcal{F}^{\pi_0}=\mathcal{M}_{\pi_0})$, cela implique régulier, cf. 4.8.5.}). Si on restreint la fcga $(\mathcal{F}^{\pi})_{\pi\geqslant\pi_0}$ aux cardinaux $\pi\geqslant\pi'_0$, où $\pi'_0\geqslant\pi_1$, alors on trouve une fcga où la majoration de cardinaux (pour $\pi=\pi^{\pi_0}$) est valable sans condition préalable $\pi'^{\pi}\geqslant\pi'_1$ ; plus précisément, on a alors}
\[
\pi'_1 \;\leqslant\; 2^{({\pi'_0}^{\pi_0})} \qquad \text{\textit{(a fortiori $\pi'_1\leqslant 2^{(2^{\pi'_0})}$)},} \tag*{(4.7.18.1)}
\]
\textit{donc la condition préalable $\pi'^{\pi}\geqslant\pi'_1$ est automatiquement vérifiée si $\pi\geqslant\pi_0$, $\pi=\pi^{\pi'_0}$, $\pi'>\pi$ (cf. remarque 4.7.7.2).}

Ceci équivaut donc à la majoration de cardinaux suivante (appliquée au cas de $\pi=\pi'_0$).

\textbf{Lemme 4.7.19.} \textit{Supposons $\pi_0$ régulier pour la fcga $(\mathcal{F}^{\pi})_{\pi\geqslant\pi_0}$. Alors pour tout cardinal $\pi\geqslant\pi_1$ ($\overset{\text{déf}}{=}\sup(\pi_0,\cardred\mathcal{F}^{\pi_0})$), on a}
\[
\cardred\mathcal{F}^{\pi} \;\leqslant\; 2^{(\pi^{\pi_0})} \, , \tag*{(4.7.19.1)}
\]
\textit{et si $\pi'\geqslant\pi_1$ et $F\in\Ob\mathcal{F}^{\pi'}$, on a}
\[
\cardred\mathcal{F}^{\pi}/F \;\leqslant\; \pi'^{(\pi^{\pi_0})} \qquad \text{(\footnotemark[78])} \, . \tag*{(4.7.19.2)}
\]
\footnotetext[78]{\textbf{N.B.} J'ignore si on peut remplacer dans ces majorations $\pi^{\pi_0}$ par $\pi_0$.}

\noindent\msp{page 185}\textsc{Démonstration du lemme}. Comme $\pi_0$ est régulier, posant $C=\mathcal{F}^{\pi_0}=\mathcal{M}_{\pi_0}$, on trouve $\mathcal{M}\simeq\Ind_{\pi_0}(C)$, donc on peut supposer $\mathcal{M}=\Ind_{\pi_0}(C)$, et appliquer les résultats prouvés pour les catégories de cette forme. Comme $\mathcal{F}^{\pi}\subset\mathcal{M}_\pi$, $\mathcal{F}^{\pi'}\subset\mathcal{M}_{\pi'}$, les majorations (4.7.9.1) et (4.7.9.2) sont contenues dans % [sic] (4.7.9.1) et (4.7.9.2) -> (4.7.19.1) et (4.7.19.2), the majorations of Lemme 4.7.19 are meant
\[
\left\{
\begin{array}{rcll}
\cardred\mathcal{M}_\pi & \leqslant & 2^{(\pi^{\pi_0})} & \\
\cardred\mathcal{M}_\pi/F & \leqslant & \pi'^{(\pi^{\pi_0})} & \text{si } F\in\Ob\mathcal{M}_{\pi'},\ \text{avec}\\
 & & & \pi'\geqslant\pi\geqslant\pi_1.
\end{array}
\right\}
\ \text{pour } \pi\geqslant\pi_1 \qquad \text{(\footnotemark[79])}. \tag*{$(*)$}
\]
\footnotetext[79]{Comparer avec majorations $\cardred\Filt^{\pi}\leqslant 2^{\pi}$, $\cardred\Filt^{\pi}/F\leqslant\pi'^{\pi}$, si $F\in\Ob\Filt^{\pi'}$, avec $\pi'\geqslant\pi\geqslant\pi_1$, majorations plus agréables que celles de (4.7.19). Il faudrait quand même s'assurer sur un exemple si les majorations plus sympathiques ne seraient pas valables dans le cas général.}
La première de ces majorations n'est autre que la deuxième majoration dans 4.7.7.1. Notons qu'on a bien
\[
2^{(\pi^{\pi_0})} \;\leqslant\; 2^{(2^{\pi})} \, ,
\]
ce qui résulte en effet de
\[
\pi^{\pi_0} \;\leqslant\; 2^{\pi} \, ,
\]
qu'on prouve en notant que $\pi<2^{\pi}$, d'où $\pi^{\pi_0}\leqslant(2^{\pi})^{\pi_0}=2^{\pi\cdot\pi_0}=2^{\pi}$ puisque $\pi\geqslant\pi_0$. Il reste à établir la deuxième majoration
\[
\cardred\mathcal{M}_\pi/F \;\leqslant\; \pi'^{(\pi^{\pi_0})} \qquad \text{si } F\in\Ob\mathcal{M}_{\pi'},\ \text{avec } \pi'\geqslant\pi\geqslant\pi_1 \, . \tag*{$(*)$}
\]
(On n'en a pas besoin pour 4.7.18, mais ce sera sans doute utile par ailleurs.) On utilise $\mathcal{M}_{\pi'}\subset\Filt^{\pi'^{\pi_0}}$ si $\pi'\geqslant\pi_1$ (4.7.6.2), donc si $F\in\Ob\mathcal{M}_{\pi'}$, on a

\end{document}
